\clearpage

\section {Introduction }

In this technical note, a description is given of the methods used by the 
Rutherford-Oxford group to measure the $\nu_{\mu}$-disappearance oscillation parameters
sin$^{2}$(2$\theta_{23}$) and $\Delta m^{2}_{23}$ from data acquired during the first T2K physics run (the 2010a dataset). 
The best-fit values of these oscillation parameters, and the 68\%, 90\% and 95\% confidence contours are also reported. 
The T2K 2010a dataset was acquired from January - June 2010 and corresponds to an integrated 
JPARC neutrino beam exposure of 3.23 $\times$ 10$^{19}$ POT. 

For this analysis, the ND280 group supplied the total number of ND280 CC-like events, but no energy spectrum was available. 
This information was used to make a prediction of the SuperK reconstructed energy spectrum of 1 $\mu$-like ring events; 
this is described in sections \ref{sec:superkprediction}.
%The predicted SuperK energy spectrum from GENIE and NEUT, and the results of this
%validation are given in section \ref{} (section 2.7).

The oscillation parameters were measured by comparing this predicted SuperK spectrum with the spectrum reconstructed from the 2010a data. 
This was done using a likelihood-ratio method as described in section \ref{sec:fitting}. Confidence contours were calculated using 
both the standard PDG method of ``constant $\Delta \chi^{2}$'' \cite{RevPartPhys} and the method of Feldman-Cousins \cite{FeldmanCousins}; 
these methods are described in section \ref{sec:contours}. 
Systematic uncertainties in the oscillation parameters are evaluated and incorporated into the confidence contour
as described in section \ref{sec:eval_osc_systematics}.

The results of the measurements of sin$^{2}$2$\theta_{23}$ and $\Delta m^{2}_{23}$ are given in section \ref{sec:oscfit_runI}. 
%Future prospects for the T2K $\nu_{\mu}$-disappearance analysis are discussed in section \ref{sec:prospects}, 
Planned improvements to the present analysis are discussed in section \ref{sec:future}.
A summary is given in section \ref{sec:summary}.  


